% Bagian: computability
% Bab: recursive-functions
% Subbagian: examples

\documentclass[../../../include/open-logic-section]{subfiles}

\begin{document}

\olfileid[id]{cmp}{rec}{exa}
\olsection{Contoh Fungsi Rekursif Primitif}


Kita sudah mempunyai beberapa contoh fungsi rekursif primitif, yaitu fungsi
penjumlahan dan perkalian~$\Add$ serta $\Mult$. Fungsi identitas
$\fn{id}(x) = x$ bersifat rekursif primitif karena fungsi itu hanyalah
$\Proj{1}{0}$. Fungsi konstan $\fn{const}_n(x) = n$ bersifat rekursif primitif
karena dapat didefinisikan dari $\Zero$ dan $\Succ$ dengan komposisi
berturut-turut. Hal ini berguna ketika kita hendak menggunakan konstanta dalam
definisi rekursif primitif. Sebagai contoh, fungsi $f(x) = 2 \cdot x$ dapat
diperoleh dengan komposisi dari $\fn{const}_2(x)$ dan perkalian sebagai
$f(x) = \Mult(\fn{const}_2(x), \Proj{1}{0}(x))$. Mulai sekarang kita akan
menggunakan siasat ini.

\begin{prop}
  Fungsi perpangkatan $\fn{exp}(x, y) = x^y$ bersifat rekursif primitif.
\end{prop}

\begin{proof}
  Kita dapat mendefinisikan $\fn{exp}$ secara rekursif primitif sebagai
  \begin{align*}
    \fn{exp}(x, 0) & = 1\\
    \fn{exp}(x, y+1) & = \Mult(x, \fn{exp}(x,y)).
    \intertext{Secara ketat, ini bukan definisi rekursif dari fungsi-fungsi
      rekursif primitif. Namun, bentuk resminya adalah:}
    \fn{exp}(x, 0) & = f(x)\\
    \fn{exp}(x, y+1) & = g(x, y, \fn{exp}(x,y)).
    \intertext{dengan}
    f(x) & = \Succ(\Zero(x)) = 1\\
    g(x, y, z) & = \Mult(\Proj{3}{0}(x, y, z), \Proj{3}{2}(x, y, z)) = x \cdot z
  \end{align*}
  sehingga $f$ dan $g$ didefinisikan dari fungsi-fungsi rekursif primitif
  dengan komposisi.
\end{proof}

\begin{prop}
  Fungsi pendahulu $\fn{pred}(y)$ yang didefinisikan oleh
  \[
  \fn{pred}(y) = \begin{cases}
    0 & \text{jika $y=0$}\\
    y-1 & \text{selain itu}
  \end{cases}
  \]
  bersifat rekursif primitif.
\end{prop}

\begin{proof}
 Perhatikan bahwa
 \begin{align*}
   \fn{pred}(0) & = 0 \text{ dan}\\
   \fn{pred}(y+1) & = y.
 \end{align*}
 Ini hampir merupakan definisi rekursif primitif. Secara ketat, definisi ini
 tidak sesuai dengan pola definisi melalui rekursi primitif karena pola itu
 mensyaratkan sekurang-kurangnya satu argumen tambahan~$x$. Definisi ini juga
 ganjil karena sama sekali tidak menggunakan $\fn{pred}(y)$ dalam definisi
 $\fn{pred}(y+1)$. Namun, mula-mula kita dapat mendefinisikan
 $\fn{pred}'(x, y)$ dengan
 \begin{align*}
   \fn{pred}'(x, 0) & = \Zero(x) = 0,\\
   \fn{pred}'(x, y+1) & = \Proj{3}{1}(x, y, \fn{pred'}(x, y)) = y.
 \end{align*}
kemudian mendefinisikan $\fn{pred}$ darinya dengan komposisi, misalnya sebagai
$\fn{pred}(x) = \fn{pred}'(\Zero(x), \Proj{1}{0}(x))$.
\end{proof}

\begin{prop}
  Fungsi faktorial $\fn{fac}(x) = \fact{x} = 1 \cdot 2 \cdot 3
  \cdot \dots \cdot x$ bersifat rekursif primitif.
\end{prop}

\begin{proof}
  Definisi rekursif primitif yang langsung adalah
  \begin{align*}
    \fn{fac}(0) &= 1\\
    \fn{fac}(y+1) & = \fn{fac}(y) \cdot (y+1).
    \intertext{Secara resmi, pertama-tama kita harus mendefinisikan fungsi
      berargumen dua $h$}
    h(x, 0) & = \fn{const}_1(x)\\
    h(x, y+1) & = g(x, y, h(x, y))
    \intertext{dengan $g(x, y, z) = \Mult(\Proj{3}{2}(x, y, z),
      \Succ(\Proj{3}{1}(x, y, z)))$, lalu menetapkan}
    \fn{fac}(y) & = h(\Proj{1}{0}(y), \Proj{1}{0}(y)) = h(y,y).
  \end{align*}
  Mulai sekarang kita akan sedikit lebih longgar dan tidak selalu memberikan
  definisi resmi melalui komposisi dan rekursi primitif.
\end{proof}

\begin{prop}
  Pengurangan terpangkas, $x \tsub y$, yang didefinisikan oleh
  \[
  x \tsub y = \begin{cases}
    0 & \text{jika $x < y$}\\
    x-y & \text{selain itu}
  \end{cases}
  \]
  bersifat rekursif primitif.
\end{prop}

\begin{proof}
  Kita mempunyai:
  \begin{align*}
    x \tsub 0 & = x\\
    x \tsub (y+1) & = \fn{pred}(x \tsub y)
  \end{align*}
\end{proof}

\begin{prop}
 Jarak antara $x$ dan $y$, yaitu $\left|x-y\right|$, bersifat rekursif primitif.
\end{prop}

\begin{proof}
  Kita mempunyai $\left| x-y \right| = (x \tsub y) + (y \tsub x)$, sehingga
  jarak dapat didefinisikan dengan komposisi dari $+$ dan $\tsub$, yang
  keduanya bersifat rekursif primitif.
\end{proof}

\begin{prop}
  Maksimum dari $x$ dan $y$, yaitu $\fn{max}(x,y)$, bersifat rekursif primitif.
\end{prop}

\begin{proof}
  Kita dapat mendefinisikan $\fn{max}(x,y)$ dengan komposisi dari $+$ dan
  $\tsub$ sebagai
  \[
  \fn{max}(x,y) \defis x + (y \tsub x).
  \]
  Jika $x$ merupakan maksimum, yakni $x \ge y$, maka $y \tsub x = 0$, sehingga
  $x + (y \tsub x) = x + 0 = x$. Jika $y$ merupakan maksimum, maka
  $y \tsub x = y - x$, sehingga $x + (y \tsub x) = x + (y - x) = y$.
\end{proof}

\begin{prop}
  \ollabel{prop:min-pr}
  Minimum dari $x$ dan $y$, yaitu $\fn{min}(x,y)$, bersifat rekursif primitif.
\end{prop}

\begin{proof}
  Latihan.
\end{proof}

\begin{prob}
  Buktikan \olref[cmp][rec][exa]{prop:min-pr}.
\end{prob}

\begin{prob}
Tunjukkan bahwa \[
f(x, y) =
2^{(2^{\iddots^{2^{x}}})}\raisebox{1ex}{\bigg\rbrace}
\raisebox{1ex}{\text {$y$ buah $2$}}\] bersifat rekursif primitif.
\end{prob}

\begin{prob}
Tunjukkan bahwa pembagian bilangan bulat $d(x, y) = \lfloor x/y \rfloor$
(yakni pembagian dengan mengabaikan seluruh bagian di belakang tanda desimal)
bersifat rekursif primitif. Jika $y = 0$, kita menetapkan $d(x, y) = 0$.
Berikan definisi eksplisit bagi~$d$ dengan rekursi primitif dan komposisi.
\end{prob}

\begin{prop}
Himpunan fungsi rekursif primitif tertutup terhadap dua operasi berikut:
\begin{enumerate}
\item Jumlah berhingga: jika $f(\vec x, z)$ bersifat rekursif primitif, maka
fungsi
\[
g(\vec x, y) \defis \sum_{z = 0}^y f(\vec x, z)
\]
juga bersifat rekursif primitif.
\item Hasil kali berhingga: jika $f(\vec x, z)$ bersifat rekursif primitif,
maka fungsi
\[
h(\vec x, y) \defis \prod_{z = 0}^y f(\vec x, z)
\]
juga bersifat rekursif primitif.
\end{enumerate}
\end{prop}

\begin{proof}
Sebagai contoh, jumlah berhingga didefinisikan secara rekursif oleh persamaan
\begin{align*}
  g(\vec x, 0) & = f(\vec x, 0)\\
  g(\vec x, y+1) & = g(\vec x, y) + f(\vec x, y+1).
\end{align*}
\end{proof}


\end{document}

